\section{The arithmetic coefficient and the limiting formula}
\label{sec:limits}

The preceding sections reduce the count of values to distinct tuples with
a common long prime prefix.  We now sum over the largest prime and replace
the remaining prefix by a volume.  The shorter tail stays discrete: its
different witnesses are combined by a union before taking that volume.
Throughout this section, $H$ is fixed first, $x$ tends to infinity next,
and $H$ tends to infinity last.  In particular, $P=\lfloor\log\log H\rfloor$,
$R=m-H$, and $L=m-P$ always refer to the two cuts in \eqref{eq:cuts}.

\subsection{Summing over the largest prime}

Let $\mathcal D_H(x)$ be the set of distinct data
$(d,p_1,\ldots,p_R)$ for which there are primes
$p_{R+1},\ldots,p_L$ and a positive integer $a$ satisfying
the basic conditions \eqref{eq:basic} with $p_0$ omitted, and
\[
 d=\varphi(w),\qquad w=p_{R+1}\cdots p_La.
\]
The convention $u_0=B$ is retained in these conditions.  Thus the
inequality with index zero is a condition on the displayed data and
their witness, independent of the eventual choice of $p_0$.
For any function $f:[1,\infty)\to[0,1]$, define
\begin{equation}
 M_f(x;H)=
 \sum_{(d,p_1,\ldots,p_R)\in\mathcal D_H(x)}
 \frac{f(\ell(d)/d)}{d\prod_{i=1}^R(p_i-1)}.
 \label{eq:mass}
\end{equation}
Each datum is included once, regardless of its number of witnesses.
The function $f=1$ makes the least-preimage ratio irrelevant to this
definition.

We use $K_H=\exp(CP\rho^{-P})=H^{o(1)}$, allowing the absolute constant
$C$ to increase in upper bounds.  Summing over witnesses, the totient
inequality
\begin{equation}
 \frac{1}{\varphi(p_{R+1}\cdots p_La)}
 \le \frac{1}{\varphi(a)}
       \prod_{i=R+1}^L\frac{1}{p_i-1}
 \label{eq:tail-reciprocal-majorant}
\end{equation}
and Lemma~\ref{lem:boxes} give
\begin{equation}
 0\le M_f(x;H)\le M_1(x;H)\ll K_HG_m.
 \label{eq:mass-upper}
\end{equation}
Indeed, the reciprocal sum for the full $L$ prime coordinates is
$O(G_m)$ by \eqref{eq:projected-mass}, while the unrestricted reciprocal
sum for the allowed smooth cofactor is $O(K_H)$ by
\eqref{eq:smooth-cost}.  The inequality in
\eqref{eq:tail-reciprocal-majorant}, rather than an equality, also
covers the possibility that $p_L$ divides $a$.

\begin{proposition}[The unweighted prime sum]
\label{prop:mass}
For each fixed $c>1$, both $t=x$ and $t=x/c$ satisfy
\begin{equation}
 V(t)=\frac{t}{\log x}
 \bigl(M_1(x;H)+E_H(t;x)G_m\bigr),
 \qquad
 \lim_{H\to\infty}\limsup_{x\to\infty}
 \max_{t\in\{x,x/c\}}|E_H(t;x)|=0.
 \label{eq:mass-count}
\end{equation}
All parameters in the two formulas are formed from the same $x$.
\end{proposition}

\begin{proof}
For fixed data in $\mathcal D_H(x)$ put
\[
 D=d\prod_{i=1}^R(p_i-1).
\]
Since $d\le w$, the deterministic bound \eqref{eq:cofactor-size}
implies
\[
 1\le D\le p_1\cdots p_La\le\exp((\log x)^{0.8}).
\]
The possible largest primes are precisely
$x^{0.9}\le p_0\le 1+t/D$.  The remaining prime factors are smaller
than $x^{0.9}$ for large $x$, so this choice introduces neither a
repeated prefix prime nor another ordering condition.  The endpoints
$t/D$ lie uniformly in $[x^{1-o(1)}/c,x]$.  The prime number theorem,
with $\log(t/D)=\log x+o(\log x)$ uniformly, therefore gives
\[
 \#\{p_0:x^{0.9}\le p_0\le1+t/D\}
 =\frac{t}{D\log x}\bigl(1+o_{x\to\infty;H,c}(1)\bigr).
\]
The lower cutoff costs a relative
$O(x^{-0.1+o(1)})$, uniformly in the data.  Summing this formula and
using \eqref{eq:mass-upper} counts $\mathcal T(t)$ with main term
$tM_1/\log x$ and error $o_{x\to\infty;H,c}(xG_m/\log x)$.
Proposition~\ref{prop:unique-prefix} replaces $\#\mathcal T(t)$ by
$V(t)$ with an error whose normalized iterated limit is zero.  Since
$t$ is either $x$ or $x/c$, this proves \eqref{eq:mass-count}.
\end{proof}

The common mass in Proposition~\ref{prop:mass} already gives the
fixed-scaling conclusion. Its proof does not require identifying that
mass with the arithmetic coefficient.

\begin{corollary}[Fixed scaling]\label{cor:fixed-scaling}
For every fixed real \(c>0\),
\[
 \frac{V(cx)}{V(x)}\longrightarrow c.
\]
\end{corollary}

\begin{proof}
First let \(c>1\). Apply Proposition~\ref{prop:mass} at \(t=x\)
and \(t=x/c\), using the same \(M_1(x;H)\) and \(G_m(x)\).
Subtracting gives
\[
 \frac{(V(x)-cV(x/c))\log x}{xG_m}
 =E_H(x;x)-E_H(x/c;x).
\]
Ford's estimate \eqref{eq:ford-scale} supplies a positive constant
\(c_-\) with \(V(x)\log x/(xG_m)\ge c_-\) for large \(x\).
Consequently
\[
 \limsup_{x\to\infty}
 \left|1-\frac{cV(x/c)}{V(x)}\right|
 \le\frac1{c_-}\limsup_{x\to\infty}
       \bigl(|E_H(x;x)|+|E_H(x/c;x)|\bigr).
\]
Let \(H\) tend to infinity and then replace \(x\) by \(cx\) to
obtain the assertion. This also covers arguments at which
\(m(x)\) and \(m(x/c)\) differ: the latter integer never enters
the comparison. Neither continuity of the phase coefficient nor
matching values at the ends of the phase interval is required.

The case \(c=1\) is immediate. If \(0<c<1\), apply the result
just proved to the multiplier \(1/c>1\) and the argument
\(y=cx\); taking reciprocals gives the stated limit.
\end{proof}

\subsection{A finite tail and a continuous prefix}

The phase-dependent witness set in \eqref{eq:AH} changes with the
phase, but for fixed $H$ all its members belong to one finite set.
This observation will permit uniform limits without any continuity
assumption.

\begin{lemma}[Uniformly finite tail data]
\label{lem:finite-tail-data}
For each sufficiently large fixed $H$, the union of all
$\mathcal W_{H,s}(d)$, over $0\le s<1$ and all $d$, is finite.
All its prime and integer coordinates, and all its totient values,
have explicit bounds depending only on $H$.  Uniformly over these
data,
\begin{equation}
 0\le D(h,\eta)\ll_H\log(2h)\qquad(h\ge H).
 \label{eq:tail-D-bound}
\end{equation}
\end{lemma}

\begin{proof}
Write $a_*=\lambda/\rho$, so that
$\lambda\le\alpha_s<a_*$.  Let
\[
 q_H=\left\lceil\exp\exp(1.1a_*H\rho^{-H})\right\rceil,
 \qquad
 A_H^*=\left\lceil\exp\exp(2a_*P\rho^{-P})\right\rceil.
\]
Every witness belongs to the finite set
$\{1,\ldots,q_H\}^{H-P}\times\{1,\ldots,A_H^*\}$.
Its corresponding integer $w$ and totient $d$ satisfy
$d\le w\le A_H^*q_H^{H-P}$.  These bounds do not involve $x$ or $s$.
Finally,
$a_j=\int_j^{j+1}\log t\,\dd t\le\log(j+1)$, so the definition
\eqref{eq:tail-D} and the bounded tail coordinates give
\eqref{eq:tail-D-bound}.
\end{proof}

For $s=\theta(x)$ and $\eta\in\mathcal W_{H,s}(d)$, define
the region $\mathcal K_\eta\subset\mathbb R^R$ by $u_0=B$ and
\begin{equation}
 u_i-\sum_{r=i+1}^Ra_{r-i}u_r\ge D(m-i,\eta)
 \qquad(0\le i\le R).
 \label{eq:prefix-region}
\end{equation}
These inequalities imply $u_i\ge0$.  The coordinates $u_i$ here are
real; the coarse prime bands have been removed.  We write
$\operatorname{Vol}_R$ for $R$-dimensional Lebesgue measure.

\begin{lemma}[The volume of the witness union]
\label{lem:tail-volume}
There are numbers $\delta_H\to0$ such that, uniformly for
$f:[1,\infty)\to[0,1]$,
\begin{equation}
 \limsup_{x\to\infty}
 \left|\frac{M_f(x;H)}{G_m}
 -\frac1{G_m}\sum_{d:\mathcal W_{H,\theta}(d)\ne\varnothing}
 \frac{f(\ell(d)/d)}d
 \operatorname{Vol}_R
 \bigcup_{\eta\in\mathcal W_{H,\theta}(d)}\mathcal K_\eta
 \right|\le\delta_H.
 \label{eq:mass-volume}
\end{equation}
Moreover, for every nonempty finite
$T\subseteq\mathcal W_{H,\theta}(d)$,
\begin{equation}
 \operatorname{Vol}_R\bigcap_{\eta\in T}\mathcal K_\eta
 =G_R\left(1-\frac1B\sum_{h=H}^m g_{m-h}
       \max_{\eta\in T}D(h,\eta)\right)_+^R.
 \label{eq:intersection-volume}
\end{equation}
Here $z_+=\max(z,0)$.
\end{lemma}

\begin{proof}
An intersection imposes the largest of the witness thresholds at each
slack coordinate.

\paragraph{Intersection volumes.}
Set $C_i=\max_{\eta\in T}D(m-i,\eta)$ and introduce the prefix
slacks
\[
 t_i=u_i-\sum_{r=i+1}^Ra_{r-i}u_r\qquad(0\le i\le R).
\]
Lemma~\ref{lem:simplex} gives
$\sum_{i=0}^Rg_it_i=B$, and the substitution from
$(u_1,\ldots,u_R)$ to $(t_1,\ldots,t_R)$ has determinant one.
The intersection is $t_i\ge C_i$ for every $i$, including $i=0$.
Translation by $C_i$ leaves a simplex of available size
$B-\sum_i g_iC_i$. Its volume is
\[
 \frac{(B-\sum_{i=0}^Rg_iC_i)_+^R}
 {R!\prod_{i=1}^Rg_i},
\]
which proves \eqref{eq:intersection-volume} after putting $h=m-i$.
It remains to compare the volume of the union with the prime mass.

\paragraph{From primes to volume.}
Identify $Q_h=p_{m-h}$ for $P\le h<H$.  These are exactly the
primes with indices $R+1,\ldots,L$.  For fixed tail witness $\eta$,
the original prefix conditions are its coarse bands together with
\begin{equation}
 \xi_i u_i-\sum_{r=i+1}^Ra_{r-i}u_r\ge D(m-i,\eta)
 \qquad(0\le i\le R).
 \label{eq:perturbed-prefix-region}
\end{equation}
Taking the union over $\eta\in\mathcal W_{H,\theta}(d)$ is essential:
a prime prefix with several witnesses contributes just once to
\eqref{eq:mass}.

First replace the reciprocal prime sum by the volume of this union
with its bands and its $\xi_i$.  In a unit box wholly contained in
the union, the product of the prime reciprocal
sums differs from its volume by a relative
$O(\sum_{i=1}^R e^{-c b_i})=O(e^{-c' b_R})$, by
Lemma~\ref{lem:boxes}.  The geometric growth of $b_i$ as $i$
decreases makes this estimate independent of $R$.  Both the total
mass and the total volume before this replacement are
$O(K_HG_m)$, by summing over witnesses as in
\eqref{eq:tail-reciprocal-majorant}.

If a unit box meets the union without being contained in it,
choose a witness whose region meets the box.  That witness's
region does not contain the box.  Consequently the box lies in
a shell at a prefix band endpoint, or within $O(h_i^2)$ of equality
in one of \eqref{eq:perturbed-prefix-region}.  For an upper bound,
sum these shells over witnesses using
\eqref{eq:tail-reciprocal-majorant}.  Thicken the discrete tail
coordinates into unit boxes as well and sum $1/\varphi(a)$ at a
cost $O(K_H)$.  The resulting full $L$-coordinate boxes are covered
by the shell bounds of Lemma~\ref{lem:boxes} and
Proposition~\ref{prop:discard}.  Their contribution is $o_H(1)G_m$.
More explicitly, the slice coefficient for an index with $h=m-i$
is $Lg_i/B\ll\rho^h$.  Summing the shell bounds therefore costs
at most
\[
 O\!\left(K_H\sum_{h\ge H}
       (h e^{-h/40}+h^3\rho^h)\right)G_m.
\]
The broader $O(K_HH^{-2})G_m$ shell estimate in
Proposition~\ref{prop:discard} also suffices.

\paragraph{Removing the slack perturbation.}
Replace $\xi_i$ by $1$ for $i\le R$, retaining the bands.
The removed region lies in a shell of thickness
$O(e^{-h_i/40}b_i)$, with the same box-thickening errors.
Its normalized contribution is at most
\[
 O\!\left(K_H\sum_{h\ge H}
       (h e^{-h/40}+h^3\rho^h)\right)=o_H(1).
\]
Both shell estimates are valid for the union because upper bounds
may sum over witnesses, while $0\le f\le1$.

\paragraph{Removing the prefix bands.}
For every point satisfying
the now unperturbed prefix inequalities, recurrence iteration gives
\begin{equation}
 u_i\ge g_{R+1-i}v_{H-1}\gg H\rho^{-h_i}
 \qquad(0\le i\le R),
 \label{eq:unbanded-prefix-lower}
\end{equation}
where $v_{H-1}=\log_2Q_{H-1}$.  To see this, retain the contribution
of the first tail coordinate in each prefix inequality, discard
the other nonnegative tail contributions, and apply the recurrence
for $g$.  The tail band for $v_{H-1}$ and
\eqref{eq:g-asymptotic} give the displayed lower bound.
Thickening the $H-P$ tail coordinates changes inequality $i$ by
$O(H\log(h_i+1))$. The enlargement term
$C h_i^2\rho^{h_i}u_i$ is at least a constant times $CHh_i^2$,
so it absorbs this error. Together with the retained tail bands,
this places the full region inside the enlarged simplex of
Lemma~\ref{lem:boxes}.

To sum these regions without a factor for the number of witnesses,
use \eqref{eq:tail-reciprocal-majorant} and first fix $a$. Partition
the tail coordinates into half-open unit boxes
$\mathcal C=\prod_{h=P}^{H-1}\mathcal C_h$. Let
$\mathcal U_{\mathcal C}$ be the union of $\mathcal K_\eta$ over
all allowed witnesses $\eta=((Q_h),a)$ with $(\log_2Q_h)_h\in\mathcal C$.
Let $E$ be the set of prefix points violating at least one removed
band. Aggregating the reciprocal weights within a box gives
\begin{align*}
 &\sum_{\substack{\eta=((Q_h),a)\text{ allowed}\\
                   (\log_2Q_h)_h\in\mathcal C}}
 \frac{\operatorname{Vol}_R(\mathcal K_\eta\cap E)}
      {\prod_h(Q_h-1)}\\
 &\quad\le
 \left(\prod_{h=P}^{H-1}
       \sum_{\log_2q\in\mathcal C_h}\frac1{q-1}\right)
       \operatorname{Vol}_R(\mathcal U_{\mathcal C}\cap E)\\
 &\quad\ll\operatorname{Vol}_L
       \bigl((\mathcal U_{\mathcal C}\cap E)\times\mathcal C\bigr).
\end{align*}
The sums over $q$ are over primes. The product is $O(1)$ uniformly
in the number of tail coordinates, by \eqref{eq:mertens-box} and
their geometrically growing lower endpoints. Here $R+(H-P)=L$
and $\mathcal C$ has unit volume.
By the preceding enlargement argument, every set on the right lies
in the same enlarged $L$-simplex. Since the tail boxes are disjoint,
summing their volumes counts a subset of that simplex once. Finally,
summing $1/\varphi(a)$ over the common smooth envelope costs $O(K_H)$.
All these bounds hold at each phase; they do not compare witness sets
at nearby phases.

Scale that enlarged simplex into the true $L$-coordinate simplex.
Its Jacobian cost is bounded, and its coordinate scaling factors
through index $R$ are
\[
 1+O\!\left(\sum_{h\ge H}
       (e^{-h/40}+h^2\rho^h)\right)=1+o_H(1).
\]
For large $H$, and then large $x$, a violation of a removed
$[0.9\widetilde b_i,1.1\widetilde b_i]$ band therefore becomes a
violation of the
$[0.95b_i,1.05b_i]$ band in the true simplex.  Lemma~\ref{lem:concentration}
applies with dimension $L=m-P$, since $h_i\ge H$ and $P/H\to0$.
Summing its bounds gives $O(e^{-cH})$ of the simplex volume.
After the smooth-cofactor factor, the additional volume is at most
$O(K_He^{-cH})G_m=o_H(1)G_m$.
This proves \eqref{eq:mass-volume}.  All constants used in these
estimates are uniform in $\theta$, since
$\lambda\le\alpha_\theta<\lambda/\rho$ and
$\alpha/\alpha_\theta=1+o_{x\to\infty;H}(1)$.

\end{proof}

\subsection{The arithmetic approximation}

\begin{proposition}[Approximation by the finite coefficient]
\label{prop:tail-limit}
There are numbers $\varepsilon_H\to0$ such that, for every
$f:[1,\infty)\to[0,1]$,
\begin{equation}
 \limsup_{x\to\infty}
 \left|\frac{M_f(x;H)}{G_m}-A_H(f;\theta(x))\right|
 \le\varepsilon_H.
 \label{eq:uniform-approximation}
\end{equation}
The bound is independent of $f$ and of phase.  Each $A_H(f;s)$ is a finite sum involving values of $f$ at explicitly
bounded rational arguments and absolutely convergent series in its
exponential weights, and $A_H(f;s)\ge0$.
\end{proposition}

\begin{proof}
Apply finite inclusion--exclusion to the union in
\eqref{eq:mass-volume}, and then use \eqref{eq:intersection-volume}.
For a nonempty subset $T$ of witnesses, write
\[
 C_T(h)=\max_{\eta\in T}D(h,\eta),\qquad
 S_T=\sum_{h=H}^{\infty}\rho^hC_T(h),\qquad
 q_{m,T}=\frac1B\sum_{h=H}^mg_{m-h}C_T(h).
\]
The series for $S_T$ converges absolutely by
\eqref{eq:tail-D-bound}.  The finite universe in
Lemma~\ref{lem:finite-tail-data} gives a uniform bound on the number
of possible witnesses and their subsets for fixed $H$; it also
makes all the following estimates uniform in phase, even when
membership in a witness set changes.

Using $B=\alpha m\rho^{-m}$ and \eqref{eq:g-asymptotic}, we obtain
\begin{align*}
 q_{m,T}
 &=\frac{\gamma}{\alpha m}
       \sum_{h=H}^m\rho^h C_T(h)
       +O_H\!\left(\frac{m\log(2m)}B\right)\\
 &=\frac{\gamma}{\alpha m}S_T+o_{x\to\infty;H}(m^{-1}).
\end{align*}
The omitted geometric tail is $O_H(\rho^m\log(2m))$.
In particular, $q_{m,T}=O_H(1/m)$, and hence
\[
 (1-q_{m,T})_+^{m-H}
 =\exp\!\left(-\frac\gamma{\alpha_\theta}S_T\right)
       +o_{x\to\infty;H}(1).
\]
Here we used the uniform scale relation
$\alpha/\alpha_\theta=1+O(\log m/m)$ and
$(m-H)q_{m,T}^2=O_H(1/m)$.  The prefactor satisfies
\begin{equation}
 \frac{G_{m-H}}{G_m}
 =\prod_{l=0}^{H-1}\frac{(m-l)g_{m-l}}B
 =\left(\frac\gamma{\alpha_\theta}\right)^H
    \rho^{H(H-1)/2}+o_{x\to\infty;H}(1).
 \label{eq:volume-prefactor-limit}
\end{equation}
There are only boundedly many terms for fixed $H$.  Thus the
inclusion--exclusion expression differs from $A_H(f;\theta(x))$
in \eqref{eq:AH} by $o_{x\to\infty;H}(1)$, uniformly in phase;
no witness multiplicity has been introduced.
Together with Lemma~\ref{lem:tail-volume}, this proves
\eqref{eq:uniform-approximation}.

Nonnegativity follows from the finite-union probability interpretation
in \eqref{eq:tail-events}: both the union probabilities and the outer
weights are nonnegative.
\end{proof}

\subsection{Uniform convergence and the asymptotic equivalent}

The preceding approximation holds for every bounded nonnegative
weight.  Convergence as $H$ grows requires a comparison quantity
that does not depend on $H$.  For $f=1$ this quantity is the
normalized count of totients.  Section~\ref{sec:companion} will
supply a different comparison quantity for each $f_k$.

\begin{lemma}[Convergence from a common comparison]
\label{lem:coefficient-convergence}
Fix $f:[1,\infty)\to[0,1]$.  Suppose a real function $C_f(x)$,
independent of $H$, and numbers $\eta_H\to0$ satisfy
\begin{equation}
 \limsup_{x\to\infty}
 \left|C_f(x)-\frac{M_f(x;H)}{G_m}\right|\le\eta_H.
 \label{eq:common-comparison}
\end{equation}
Then $A_H(f;s)$ converges uniformly for $0\le s<1$ to the
limit $A(f;s)$ defined in Section~\ref{sec:statements}, and
\begin{equation}
 C_f(x)=A(f;\theta(x))+o(1).
 \label{eq:comparison-limit}
\end{equation}
If additionally $c_-\le C_f(x)\le c_+$ for all sufficiently large
$x$, then $c_-\le A(f;s)\le c_+$ for every $s\in[0,1)$.
\end{lemma}

\begin{proof}
By Proposition~\ref{prop:tail-limit},
\[
 \limsup_{x\to\infty}
 |C_f(x)-A_H(f;\theta(x))|\le\zeta_H,
 \qquad \zeta_H=\eta_H+\varepsilon_H\longrightarrow0.
\]
Every phase occurs at arbitrarily large real arguments.  Indeed,
\[
 \psi(B)=\frac{\log B-\log\log B}{\lambda}
\]
is continuous and strictly increasing for $B>e$, and tends to
infinity.  For any $s\in[0,1)$ solve $\psi(B_n)=n+s$ and put
$x_n=\exp(\exp B_n)$.  Then $m(x_n)=n$ and $\theta(x_n)=s$
exactly.  Comparing indices $H$ and $K$ along this same sequence
gives
\[
 |A_H(f;s)-A_K(f;s)|\le\zeta_H+\zeta_K.
\]
This bound is independent of $s$, so the functions form a uniformly
Cauchy sequence.  In particular,
$\sup_s|A_H(f;s)-A(f;s)|\le\zeta_H$.
Combining this inequality with the preceding limsup bound and
then sending $H$ to infinity proves \eqref{eq:comparison-limit}.
If $C_f$ lies between $c_-$ and $c_+$, apply those bounds on each
exact-phase sequence and let $H$ grow to obtain the asserted bounds
for $A$.  No continuity in the phase was used.
\end{proof}

\begin{proof}[Proof of Theorem~\ref{thm:main}]
Set
\[
 C_1(x)=\frac{V(x)\log x}{xG_m}.
\]
Proposition~\ref{prop:mass} with $t=x$ supplies
\eqref{eq:common-comparison}.  Lemma~\ref{lem:coefficient-convergence}
therefore proves uniform existence of $A(1;s)$ and
$C_1(x)=A(1;\theta(x))+o(1)$.  Ford's two-sided estimate
\eqref{eq:ford-scale} supplies positive absolute constants
$c_-,c_+$ bounding $C_1$.  The same lemma gives
\[
 0<c_-\le A(1;s)\le c_+<\infty\qquad(0\le s<1).
\]
The additive approximation is consequently equivalent to
\[
 V(x)\sim\frac{x}{\log x}G_mA(1;\theta(x)).
\]
The coefficient in this formula was defined by the finite arithmetic
expressions \eqref{eq:AH}; the count $V$ has been used to establish
their convergence and bounds, not to define them.

Corollary~\ref{cor:fixed-scaling} supplies the remaining assertion of
the theorem for every \(c>0\).
\end{proof}
